\section{Introduction}
\label{sec:intro}

%% --- 1. Context -------------------------------------------------
Large language model (LLM) agents now run as sequential software
systems in production: a deployed agent plans, calls tools, observes
results, retries, and terminates in success or failure, as in
ReAct-style and tool-using
designs~\cite{yao2023react,schick2023toolformer,wang2024agentsurvey}.
Once such an agent sits on a critical path---resolving a support
ticket, patching a repository, issuing a refund---the operator inherits
the ordinary obligations of dependability and site reliability
engineering (SRE)~\cite{avizienis2004dependability,beyer2016sre}:
state a reliability target, quantify the uncertainty around it, and say
under what conditions it holds.

%% --- 2. Narrow to the specific problem --------------------------
Agent evaluation, however, reports reliability as a scalar. Benchmarks
supply rich execution traces---ToolBench~\cite{qin2023toolllm},
$\tau$-bench~\cite{taub2024}, SWE-bench~\cite{swebench2024},
WebArena~\cite{zhou2023webarena}, and
AgentBench~\cite{agentbench2024}---but summarize them as
pass$@k$~\cite{chen2021codex}, pass$^k$, or the reliability decay curve
(RDC)~\cite{rdc2025}. Each score is a different summary of when a run
succeeds, computed with a different denominator and reported without a
fit check. Classical software reliability already has the machinery for
this setting, from Markov program models~\cite{cheung1980} and
absorbing chains~\cite{kemenysnell} to reliability-growth
models~\cite{goelokt,musa1987}, and trace-driven frameworks such as
TriCEGAR~\cite{tricegar2026} and ProbGuard~\cite{pro2guard2025} learn
Markov abstractions from execution data.

%% --- 3. The gap -------------------------------------------------
Despite this, no current LLM-agent reliability report names the
success-time distribution it estimates, tests whether the traces
support that distribution, or attaches finite-trace uncertainty to it.
The practical cost is concrete. A team that ships a refund agent at
$\mathrm{pass}@1 = 0.72$ cannot say what reliability a larger step
budget buys, what an added fallback tool changes end-to-end, or how an
internal $\mathrm{pass}^5$ dashboard relates to an SRE
mean-time-between-failures target---without rerunning the benchmark
once per question, and still without an interval or a validity check.

%% --- 4. Contribution --------------------------------------------
In this paper we propose \textsc{TraceToChain}, a pipeline that turns an
agent trace corpus into an \emph{audited} reliability model. Each
execution becomes a trajectory in an absorbing discrete-time Markov
chain (DTMC) $\hat M=(\ST,\hat Q,\hat R_\oplus,\hat R_\ominus,s_0)$
(Definition~\ref{def:amc}), with transient states $\ST$ for
intermediate behavior and absorbing states for terminal success and
failure. What distinguishes \textsc{TraceToChain}
from prior trace-to-DTMC abstraction is that the fit is not trusted by
default: a composite Akaike information criterion (AIC) and
Kolmogorov--Smirnov (KS) certificate gates every reported quantity, and
Dirichlet-posterior and bootstrap intervals accompany the ones that
pass. Horizon reliability, local perturbation, and metric
reconciliation then become first-passage queries on one model
(Props.~\ref{thm:closed} and~\ref{thm:perturb}, Thm.~\ref{thm:unify}),
so pass$@k$, pass$^k$, and the RDC stop competing and become
projections of a single success-time distribution.

%% --- 5. Approach and key results in brief -----------------------
Figure~\ref{fig:design} shows the pipeline and the study design that
evaluates it: a strict 50/50 fit/test protocol on seven controlled
MAST-style corpora, then end-to-end fits on $2{,}459$ real agent
episodes from two independent benchmarks. In both settings the
certified outputs track their held-out targets closely, and the
certificate marks precisely which further outputs its evidence
supports---the behavior a reliability audit is supposed to have.

%% --- 6. Explicit contributions ----------------------------------
Our contributions are:
\begin{enumerate}[(1)]
  \item \emph{An audited trace-to-chain pipeline.}
  \textsc{TraceToChain} (\S\ref{sec:state_construction},
  Algorithms~\ref{alg:trace}--\ref{alg:gof}) maps traces to a fitted
  absorbing DTMC behind a composite AIC$\,\wedge\,$KS goodness-of-fit
  certificate, with Dirichlet-posterior and trace-level bootstrap
  intervals on the chain and every derived quantity (\S\ref{sec:uq}).
  \item \emph{Metric unification.}
  pass$@k$, pass$^k$, and the RDC are derived as projections of one
  success first-passage distribution (Thm.~\ref{thm:unify}), making
  their denominators explicit and their disagreement resolvable.
  \item \emph{Held-out validation.}
  On seven controlled MAST-style frameworks (SS9, \S\ref{sec:ss9}) the
  pipeline attains maximum RDC error
  $L_\infty^{\mathrm{RDC}}=0.053$ (median $0.048$) and held-out KS
  $p>0.05$ on $7/7$ frameworks (min $p=0.83$).
  \item \emph{Real-trace validation on two benchmarks.}
  Fitting end-to-end to $479$ SWE-bench~Verified SWE-agent trajectories
  (SS10--SS14, \S\ref{sec:ss10}) and $1{,}980$ $\tau$-bench episodes
  (SS15, \S\ref{sec:ss15}) recovers the held-out pass rate to within
  $0.03$ and to a mean of $0.0075$ respectively, delivering
  metric-reconciliation and closed-form perturbation answers on real
  agents while the audit scopes the remaining outputs.
  \item \emph{Sensitivity analysis.}
  Clustering, censoring, and bootstrap-path studies
  (SS11--SS13, \S\ref{sec:robustness}) show the verdict is invariant
  across $18$ clustering configurations, and that $\hat R_\infty$ stays
  within $4\%$ of a synthetic ground truth up to $50\%$ censoring.
\end{enumerate}

%% --- 7. Roadmap -------------------------------------------------
The rest of this paper is organized as follows.
\S\ref{sec:related} positions the work; \S\ref{sec:formal} and
\S\ref{sec:state_construction} define and fit the absorbing-chain
model; \S\ref{sec:core} and \S\ref{sec:ext} derive the reliability
results; \S\ref{sec:sim}--\S\ref{sec:robustness} report the controlled,
real-trace, and sensitivity studies; and
\S\ref{sec:threats}--\S\ref{sec:conc} discuss limits and implications.

\begin{figure}[!t]
\centering
\begin{tikzpicture}[
  >=Stealth, font=\scriptsize\sffamily,
  st/.style={rectangle, rounded corners=2pt, draw=blue!65!black, fill=blue!6,
             line width=0.6pt, align=center, text width=1.45cm,
             minimum height=0.78cm, inner sep=1.5pt},
  gate/.style={diamond, aspect=1.5, draw=red!60!black, fill=red!8,
               line width=0.6pt, align=center, inner sep=0.5pt, font=\tiny\sffamily},
  good/.style={rectangle, rounded corners=2pt, draw=green!55!black, fill=green!8,
               line width=0.7pt, align=center, text width=2.55cm,
               minimum height=0.95cm, inner sep=2pt},
  bad/.style={rectangle, rounded corners=2pt, draw=orange!70!black, fill=orange!9,
              line width=0.6pt, align=center, text width=1.75cm,
              minimum height=0.95cm, inner sep=2pt},
  ev/.style={rectangle, rounded corners=2pt, draw=gray!55, fill=gray!5,
             line width=0.5pt, align=left, text width=2.28cm, inner sep=2.5pt,
             font=\tiny\sffamily},
  ar/.style={->, draw=gray!75, line width=0.7pt},
  dar/.style={->, draw=gray!50, line width=0.5pt, dashed, shorten >=1pt}]
  \node[st] (tr)  at (0.80, 0)     {agent\\traces $\mathcal T$};
  \node[st] (fe)  at (3.15, 0)     {featurize $\phi$\\+ cluster};
  \node[st] (mle) at (5.50, 0)     {Laplace MLE\\$\hat Q,\hat R_\oplus,\hat R_\ominus$};
  \node[gate] (g) at (5.50, -1.45) {AIC\\$\wedge$\,KS};
  \node[good] (ok) at (1.75, -1.45)
      {\textbf{certified}\\$\mathcal R(d)$, $R_\infty$, pass$@k$,\\pass$^k$, $\Delta R_\infty$ $+$ 95\% CI};
  \node[bad] (no) at (7.55, -1.45) {\textbf{withheld}\\re-featurize /\\richer model};
  \draw[ar] (tr) -- (fe);
  \draw[ar] (fe) -- (mle);
  \draw[ar] (mle) -- (g);
  \draw[ar] (g) -- node[above, font=\tiny] {pass} (ok);
  \draw[ar] (g) -- node[above, font=\tiny] {fail} (no);
  \draw[gray!35, line width=0.4pt] (-0.05,-2.42) -- (8.45,-2.42);
  \node[anchor=west, font=\tiny\itshape, text=gray!45!black]
       at (-0.05,-2.63) {Evaluation design};
  \node[ev] (e1) at (1.25, -3.62)
      {\textbf{Controlled}\\SS1 closed form vs.\ MC\\SS7 gate Type-I/II\\SS9 held-out, 7 corpora};
  \node[ev] (e2) at (4.20, -3.62)
      {\textbf{Real traces}\\SS10/SS14 SWE-bench\\($479$ trajectories)\\SS15 $\tau$-bench ($1{,}980$)};
  \node[ev] (e3) at (7.15, -3.62)
      {\textbf{Sensitivity}\\SS11 clustering ($18$)\\SS12 censoring ($0$--$50\%$)\\SS13 bootstrap path};
  \draw[dar] (e1.north) -- ++(0,0.42);
  \draw[dar] (e2.north) -- ++(0,0.42);
  \draw[dar] (e3.north) -- ++(0,0.42);
\end{tikzpicture}
\caption{\textsc{TraceToChain} pipeline and evaluation design.}
\label{fig:design}
\end{figure}

%% [REDUCTION] The three-column motivating figure is suppressed for the
%% 12-page body; the paragraph above states Q1--Q3 and points at the
%% results that answer them. Retained in source for the artifact.
\iffalse
\begin{figure}[!t]
\centering
\begin{tikzpicture}[
    >=Stealth,
    font=\scriptsize\sffamily,
    hdr/.style={font=\bfseries\sffamily\scriptsize, align=center},
    q/.style={rectangle, rounded corners, draw=orange!80!black,
              fill=orange!10, line width=0.7pt, align=left,
              text width=2.15cm, inner sep=2.5pt, minimum height=1.25cm},
    s/.style={rectangle, rounded corners, draw=red!60!black,
              fill=red!6, line width=0.7pt, align=left,
              text width=2.15cm, inner sep=2.5pt, minimum height=1.25cm},
    a/.style={rectangle, rounded corners, draw=green!55!black,
              fill=green!6, line width=0.7pt, align=left,
              text width=2.15cm, inner sep=2.5pt, minimum height=1.25cm},
    lbl/.style={font=\scriptsize\itshape, text=gray!50!black},
    arr/.style={->, draw=gray!70, line width=0.7pt},
    xscale=1.0, yscale=1.0,
]
    % column centers
    \def\colQ{0}
    \def\colS{2.75}
    \def\colA{5.5}
    % row y-coordinates
    \def\rowH{0.85}
    \def\rowA{-0.5}
    \def\rowB{-2.05}
    \def\rowC{-3.6}

    % Headers
    \node[hdr] at (\colQ, \rowH)
        {Question};
    \node[hdr] at (\colS, \rowH)
        {Scalar metric};
    \node[hdr] at (\colA, \rowH)
        {Fitted chain $\hat M$};

    % Row 1: budget
    \node[q] (q1) at (\colQ, \rowA)
        {\textbf{Q1.} If $d{=}8\!\to\!16$, what reliability?};
    \node[s] (s1) at (\colS, \rowA)
        {Re-run at each $d$. No interval or extrapolation.};
    \node[a] (a1) at (\colA, \rowA)
        {$\mathcal R(d)$ on any grid
         (Prop.~\ref{thm:closed}) + 95\% interval.};

    % Row 2: fallback tool
    \node[q] (q2) at (\colQ, \rowB)
        {\textbf{Q2.} Add fallback search: $\Delta R_\infty$?};
    \node[s] (s2) at (\colS, \rowB)
        {Re-run full benchmark for one new point.};
    \node[a] (a2) at (\colA, \rowB)
        {Bound $\Delta R_\infty$
         locally (Prop.~\ref{thm:perturb}).};

    % Row 3: metric zoo
    \node[q] (q3) at (\colQ, \rowC)
        {\textbf{Q3.} pass$^5$ vs.\ SRE-MTBF: which is right?};
    \node[s] (s3) at (\colS, \rowC)
        {Two denominators and ad hoc reconciliation.};
    \node[a] (a3) at (\colA, \rowC)
        {Both are marginals of one success-time distribution
         (Thm.~\ref{thm:unify}).};

    \draw[arr] (q1) -- (s1);
    \draw[arr] (s1) -- (a1);
    \draw[arr] (q2) -- (s2);
    \draw[arr] (s2) -- (a2);
    \draw[arr] (q3) -- (s3);
    \draw[arr] (s3) -- (a3);
\end{tikzpicture}
\caption{Motivating example: a fitted absorbing chain turns one trace
         corpus into horizon, perturbation, and metric-reconciliation
         queries.}
\label{fig:motivating}
\end{figure}
\fi


%% ----------------------------------------------------------------
%% [REDUCTION 2026-04-24] Expanded contributions block retained
%% for the artifact but replaced above with a prose paragraph
%% to save ~30 lines in the 12-page body.
%% ----------------------------------------------------------------
\iffalse
\begin{enumerate}[(C1)]
  \item \emph{\textsc{TraceToChain} pipeline and goodness-of-fit
        certificate} (\S\ref{sec:state_construction},
        Algorithms~\ref{alg:trace}--\ref{alg:gof}): agglomerative
        Ward clustering with silhouette-chosen $m$, Laplace-smoothed
        MLE with explicit time-homogeneity, variable-horizon, and
        right-censoring semantics, plus a composite AIC$\,\wedge\,$KS
        test that certifies whether an agent trace corpus is
        adequately modeled by a first-order absorbing DTMC.
  \item \emph{Uncertainty quantification}
        (\S\ref{sec:uq}, \texttt{mcr.uncertainty}): closed-form
        Dirichlet posterior credible intervals on every entry of
        $\hat Q, \hat R_\oplus, \hat R_\ominus$ (Beta-marginal
        Dirichlet-Multinomial conjugacy) plus a trace-level
        non-parametric bootstrap. Every downstream reliability
        quantity inherits a calibrated CI.
  \item \emph{Held-out empirical validation}
        (\S\ref{sec:ss7}, \S\ref{sec:mast}): we fit on a
        randomly-sampled half of each framework's MAST-derived
        corpus and report, on the truly held-out half, (i)~the
        empirical RDC overlaid on the analytic $\mathcal R(d)$,
        (ii)~the Kolmogorov--Smirnov distance $D_{\mathrm{KS}}$ and
        $p$-value, and (iii)~the $L_\infty$ RDC error.
        SS7 additionally validates the Type-I/Type-II error of the
        composite test on synthetic chains with known ground truth.
  \item \emph{Metric unification} (Theorem~\ref{thm:unify},
        Corollary~\ref{cor:diag}): pass$@k$, pass$^k$, and the RDC
        are all restrictions of the same first-passage distribution.
        A correlated-trial correction gives a Jensen-gap diagnostic
        that flags when independent-trial assumptions fail.
  \item \emph{Classical analytic tools, application-layer}
        (Propositions~\ref{thm:closed}--\ref{thm:mixing}):
        closed-form $\mathcal R(d)$, $R_\infty$, and $\mathrm{MTTA}$;
        a condition-number sensitivity bound; RDC-shape sufficient
        conditions; the Goel--Okumoto NHPP limit; and a spectral
        horizon bound. Each is attributed and framed as an
        application of an established result.
  \item \emph{Reproducibility artifact}: the entire pipeline (six
        simulation studies, the MAST corpus, and all tables /
        figures in this paper) runs on a laptop in under
        $15$ minutes with a single command.
\end{enumerate}
\fi
%% ----------------------------------------------------------------
